% Мэдлэгийн модны бүтэц (TikZ). Эх: config/knowledge-topics.json (2026.11.16-нд тулгасан):
% үндэс -> 11 хүрээ -> 31 дэд сэдэв. Зурагт «Элсэлт», «Байршил, өрөө» хоёр хүрээний дэд сэдвийг жишээ болгон задалсан.
% 2026.11.20 (багш нар 11/17: «өнгө бүгд ижил»): хүрээг мэдээллийн эх сурвалжаар нь өнгөөр ялгасан —
% хөх = вэб сайтаас авсан 8 хүрээ, ногоон = сургуулийн зүгээс өгсөн мэдээллийн 3 хүрээ (Хүснэгт 3.1-тэй ижил);
% дэд сэдэв нь хүрээнийхээ өнгөтэй. Доод талд тайлбартай.
% 2026.11.23 (багш: «жишээ болгон задалсан хүрээ өөр өнгөтэй байх хэрэг байна уу, бусад шиг болго»):
% задалсан хоёр хүрээний тод дүүргэлтийг авч бусадтай ижил болгов; гарчгийн «(11)» тоог хасав (тайлбар доор нь).
\begin{tikzpicture}[
    font=\footnotesize,
    node distance=0pt,
    box/.style={draw=nmBlueD, fill=nmBlue, line width=0.7pt, rounded corners=1.5pt, align=left, inner sep=3pt, minimum height=5.2mm},
    dom/.style={box, text width=4.3cm, anchor=west},
    school/.style={draw=nmGreenD, fill=nmGreen},
    subt/.style={box, text width=6.6cm, anchor=west},
    key/.style={minimum width=4mm, minimum height=3mm, inner sep=0pt, line width=0.7pt},
    note/.style={font=\scriptsize, inner sep=1.5pt},
    line/.style={semithick, black!70}]

\node[box, draw=nmVioletD, fill=nmViolet, align=center, text width=2.3cm] (root) at (0,-3.25) {Сургуулийн\\мэдээлэл};

\node[dom]       (d0)  at (2.9,0)     {Сургуулийн тухай};
\node[dom]       (d1)  at (2.9,-0.65) {Элсэлт};
\node[dom]       (d2)  at (2.9,-1.30) {Сургалт, тэнхимүүд};
\node[dom]       (d3)  at (2.9,-1.95) {Оюутны амьдрал};
\node[dom]       (d4)  at (2.9,-2.60) {Карьер, хамтын ажиллагаа};
\node[dom]       (d5)  at (2.9,-3.25) {Мэдээ};
\node[dom]       (d6)  at (2.9,-3.90) {Арга хэмжээ};
\node[dom]       (d7)  at (2.9,-4.55) {Холбоо барих};
\node[dom, school] (d8)  at (2.9,-5.20) {Байршил, өрөө};
\node[dom, school] (d9)  at (2.9,-5.85) {Багш, ажилтан};
\node[dom, school] (d10) at (2.9,-6.50) {Анги, хичээлийн хуваарь};

\foreach \n in {d0,d1,d2,d3,d4,d6,d7,d8,d9,d10}
    \draw[line] (root.east) -- ++(0.45,0) |- (\n.west);
\draw[line] (root.east) -- (d5.west);

\node[subt] (a1) at (8.6,0.35)  {Элсэх шаардлага};
\node[subt] (a2) at (8.6,-0.30) {Шалгалт, эссэ, ярилцлага};
\node[subt] (a3) at (8.6,-0.95) {Бүрдүүлэх материал};
\node[subt] (a4) at (8.6,-1.60) {Төлбөр, хураамж};
\node[subt] (a5) at (8.6,-2.25) {Бүртгэл, элсэлтийн журам};
\node[subt] (a6) at (8.6,-2.90) {Олимпиад, холбогдох элсэлтийн нөхцөл};
\foreach \n in {a1,a2,a3,a4,a5,a6}
    \draw[line] (d1.east) -- ++(0.45,0) |- (\n.west);

\node[subt, school] (l1) at (8.6,-4.55) {Давхар бүрийн өрөө};
\node[subt, school] (l2) at (8.6,-5.20) {Ресепшн, үйлчилгээний өрөө};
\node[subt, school] (l3) at (8.6,-5.85) {Тэнхимийн өрөө};
\foreach \n in {l1,l3}
    \draw[line] (d8.east) -- ++(0.45,0) |- (\n.west);
\draw[line] (d8.east) -- (l2.west);

\node[font=\scriptsize\bfseries, anchor=south] at (root.north |- 0,0.75) {Үндэс};
\node[font=\scriptsize\bfseries, anchor=south west] at (2.9,0.75) {Хүрээ};
\node[font=\scriptsize\bfseries, anchor=south west] at (8.6,0.75) {Дэд сэдэв (жишээ)};

% --- өнгөний тайлбар
\node[key, draw=nmBlueD, fill=nmBlue] (k1) at (3.1,-7.3) {};
\node[note, anchor=west] at (k1.east) {вэб сайтаас авсан мэдээлэл};
\node[key, draw=nmGreenD, fill=nmGreen] (k2) at (8.8,-7.3) {};
\node[note, anchor=west] at (k2.east) {сургуулийн зүгээс өгсөн мэдээлэл};
\end{tikzpicture}
